\section{Background}
\label{sec:background}

\subsection{Two transport protocols, and an expectation}

For all practical purposes the Internet has two transport protocols. TCP
provides a reliable ordered byte stream, and pays for it: every segment is
acknowledged, the sender maintains state to retransmit what is not, and a
congestion controller decides when to send. UDP provides none of that. It adds
a port pair and a checksum to a datagram and hands it up.

The expectation that follows is natural and, until recently, was correct: where
the application does not need what TCP provides---because the medium is
lossless, or because the application does its own recovery---UDP should be
faster, since it does strictly less work per byte.

At 100\,GbE that expectation inverts. Table~\ref{tab:opening} gives the
comparison on one receiving core of our testbed, ten runs each. A single TCP
connection delivers 56.6\,Gbps with nobody having configured anything. A single
UDP flow, on the same core and the same NIC, cannot be driven past 38.7 at any
offered rate we tried, and at 56\,Gbps offered it delivers 35.4 while discarding
37\% of what arrived.

\begin{table}[t]
\centering
\small
\begin{tabular}{@{}lrr@{}}
\toprule
 & one flow & eight flows \\
\midrule
TCP, nothing configured & \textbf{56.6} & 44.3 \\
UDP, \texttt{rmem\_default} = 1\,MB & 38.7 & 55.7 \\
UDP, \texttt{rmem\_default} = 8\,MB & 55.6 & 22.9 \\
\midrule
UDP with \sys & \textbf{62.7} & \textbf{55.8} \\
\bottomrule
\end{tabular}
\caption{Goodput in Gbps on one receiving core, $n=10$; the UDP figures are the
best any offered rate produced. TCP obtains its number by sizing its own receive
buffer. UDP has no mechanism that does so, and neither fixed value serves both
columns.}
\label{tab:opening}
\end{table}

The protocol that does less is a third slower, doing nothing the other one does.
It is worth being precise about what this is not. It is not a tuning oversight:
38.7 is the best figure any offered rate produced, not a point we failed to
optimise. And it is not uniform---the eight-flow column reverses, for reasons
\S\ref{sec:problems} makes exact---which is itself part of the problem, since it
means no single setting is right.

\subsection{What UDP did not receive}

The gap is not in the datagram path's instruction count, which remains the
shorter of the two. It is in everything that was added around TCP and not
around UDP.

TCP's receive side sizes its own buffer. \texttt{tcp\_rcv\_space\_adjust()}
measures the bytes delivered to the application over one round trip and sets
\texttt{sk\_rcvbuf} to about twice that, so the buffer tracks the
bandwidth--delay product as it changes. This was added in 2004 and has been
refined since; no operator on our testbed did anything to obtain the 56.6 above.

UDP has no equivalent, in any kernel version including the current one. A UDP
socket receives \texttt{net.core.rmem\_default} and keeps it. The admission test
in \texttt{\_\_udp\_enqueue\_schedule\_skb()} reads
\texttt{READ\_ONCE(sk->sk\_rcvbuf)} and compares; nothing anywhere adjusts that
value in response to anything.

This is not an oversight so much as a consequence of where attention went. UDP
was, for most of its life, the protocol for small request--response exchanges
and for applications willing to lose data. Neither case is sensitive to the
receive buffer. What changed is that UDP now carries QUIC, and through QUIC a
large and growing share of web traffic; it carries media transport, telemetry,
and storage replication. These are bulk transfers at line rate, and they are
sensitive to exactly the thing nobody built.

The direction of travel makes this worse rather than better. 100\,GbE is now the
baseline in the machines we measure; 400\,GbE is deployed, and 800\,GbE and
1.6\,Tbps are specified. Per-core receive cost is what binds at these rates, and
it is the quantity that has not improved.

\subsection{The receive path as it stands}

A packet arriving at the NIC is written by DMA into a page the driver posted on
a receive ring, and a completion entry is placed on a completion queue. The
driver does not take an interrupt per packet: it asks the NIC to wait until a
time or a count threshold is met (\texttt{rx-usecs} and \rxframes). On the
interrupt, NAPI polls the completion queue, builds socket buffers, and pushes
them up the stack, where GRO may coalesce consecutive packets of one flow into
one large \texttt{skb}.

For UDP, delivery means \texttt{\_\_udp\_enqueue\_schedule\_skb()}, which admits
the packet if the socket's accounted memory plus this packet does not exceed
\rcvbuf, and drops it otherwise. \texttt{recvmsg} copies out and releases the
accounting. There is no flow control: a UDP sender is not told that the
receiver's queue is full, and dropping is the receiver's only response.

Recent kernels have improved the producer side---6.18 replaced a per-socket
spinlock with a per-node lockless list, so producers batch into the receive
queue under one lock acquisition---and our results are on top of that. The
sizing question is orthogonal to the locking question and outlived it.

\subsection{Five problems}
\label{sec:problems}

What follows is what we measured to be wrong, stated as the design has to
address it. Each is established in \S\ref{sec:three} or \S\ref{sec:workingset};
\S\ref{sec:design} takes them in the same order.

\paragraph{P1: the buffer must hold a burst whose size nobody tells the
application.} The socket buffer must absorb what one moderation window delivers
while the consumer works through the previous one. On our hardware the driver's
adaptive controller settles on 128 frames, which at a 9000-byte MTU is about
1.15\,MB arriving between wake-ups. A 1\,MB buffer---a value an operator might
well pick---overflows on every window. The frame count is not exposed to the
socket layer, and the driver changes it at run time in response to load.

\paragraph{P2: what the buffer may cost is a global quantity no socket can
see.} Throughput falls off a cliff when the total the receive path touches---the
NIC's descriptor pages plus every socket's queued data---exceeds the last-level
cache. The threshold is a property of the sum. Eight sockets at 2\,MB and one at
16\,MB collapse at the same aggregate, so a per-socket limit cannot express it:
8\,MB is entirely safe for one socket and ruinous as the eighth.

\paragraph{P3: past the consumer's capacity, a larger buffer is actively
harmful.} What a large buffer costs is not the memory but the hand-off. Data
queued in a small buffer is still cache-resident when the consumer reaches it;
data queued in a large one has been evicted. Separating producer from consumer
and varying only the cache level they share, the 512\,KB-to-18\,MB throughput
ratio is 0.40 when they share L2 and L3, and 0.94 when they share nothing.
Consequently ``tune it larger'' is not merely insufficient past the ceiling, it
is the wrong direction.

\paragraph{P4: the work is spent before the drop.} A packet dropped at the
socket has already cost descriptor processing, an \texttt{skb} allocation, GRO,
and a route lookup. Under overload the receiving core spends most of itself on
packets it will discard, which is why throughput collapses with the core
measurably idle---63\% busy delivering 33\,Gbps against 56 offered.

\paragraph{P5: buffers are per socket, the bottleneck is per queue.} When TCP
and UDP share a receive queue, no amount of UDP buffer tuning gives TCP its
share, because what TCP is short of is not buffer. Allocation cannot arbitrate
between protocols; only something acting at the queue can.
